% SPDX-License-Identifier: CC-BY-SA-4.0
% Author: Matthieu Perrin

\begingroup

\begin{exercice}[Révisions -- Le lemme d'Arden]\label{exo:lexing/arden/revision}

  \begin{question}

  \item Résoudre le système d'équations rationnelles suivant :
    $\left\{
    \begin{array}{rcl}
      L_1 &=& aL_2 \mid cL_3 \\
      L_2 &=& bL_2 \mid aL_4 \mid \varepsilon \\
      L_3 &=& cL_3 \mid bL_4 \\
      L_4 &=& aL_4 \mid \varepsilon
    \end{array}
    \right.$

  \item Résoudre le système suivant :
    $\left\{
    \begin{array}{rcl}
      L_1 &=& bL_1 \mid aL_2 \\
      L_2 &=& cL_3 \mid aL_4 \\
      L_3 &=& bL_3 \mid \varepsilon \\
      L_4 &=& cL_4 \mid \varepsilon
    \end{array}
    \right.$

  \item Montrer que $a^+ \mid  a^\star  (ba^\star )^+$ est solution de $X = (a \mid b) X \mid a\mid b$.
  \item En déduire $a^+ \mid  a^\star  (ba^\star )^+ \equiv (a\mid b)^+$. 
  \end{question}

\end{exercice}

\begin{correction}

  \begin{question}

  \item
    $L_4 = a^\star$\\
    $L_2 = b^\star(aL_4 \mid \varepsilon) = b^\star(a a^\star \mid \varepsilon)= b^\star a^\star$\\
    $L_3 = c^\star bL_4 = c^\star b a^\star$\\
    $L_1 = aL_2 \mid cL_3 = ab^\star a^\star \mid cc^\star b a^\star$
 
  \item
    $L_3 = b^\star$\\
    $L_4 = c^\star$\\
    $L_2 = cL_3 \mid aL_4 = cb^\star \mid ac^\star$\\
    $L_1 = b^\star aL_2 = b^\star a(cb^\star \mid ac^\star) = b^\star acb^\star \mid b^\star aa c^\star$
 
  \item $\begin{array}[t]{rcll}
    &      & (a \mid b) (\underline{a^+ \mid  a^\star  (ba^\star )^+}) \mid a\mid b\\
    &\equiv& a a^+ \mid  a^+  (ba^\star )^+ \mid ba^+ \mid  ba^\star  (ba^\star )^+ \mid a\mid b &\text{Développement}\\
    &\equiv& a^+ \mid  a^+  (ba^\star )^+ \mid ba^+ \mid  ba^\star  (ba^\star )^+ \mid b & aa^+ \mid a \equiv a^+\\
    &\equiv& a^+ \mid  a^+  (ba^\star )^+ \mid ba^\star \mid  ba^\star  (ba^\star )^+ & ba^+ \mid b \equiv ba^\star\\
    &\equiv& a^+ \mid  a^+  (ba^\star )^+ \mid (ba^\star )^+ &  LL^+ \mid L \equiv L^+\\
    &\equiv& \underline{a^+ \mid  a^\star  (ba^\star )^+} \\
  \end{array}$
 
  \item D'après le lemme d'Arden, l'unique solution de $X = (a \mid b) X \mid a\mid b$ est $(a \mid b)^+$, \\donc $a^+ \mid  a^\star  (ba^\star )^+ \equiv (a \mid b)^+$. 

  \end{question}

\end{correction}

\endgroup
\endinput
